%=================================================
% Economics Coursework Template
%=================================================
\documentclass[12pt]{article}

%------------ Packages ------------
\usepackage[margin=1in]{geometry}   % Page layout
\usepackage{lmodern}                % Latin Modern (LaTeX default font)
\usepackage{amsmath, amssymb}       % Math symbols
\usepackage{booktabs}               % Professional tables
\usepackage{graphicx}               % Figures
\usepackage{natbib}                 % References
\usepackage[hidelinks]{hyperref}    % Hyperlinks
\usepackage{fancyhdr}               % Headers/footers
\usepackage{setspace}               % Line spacing
\usepackage{titlesec}               % Section formatting
\usepackage{comment} 
\usepackage{xcolor} 
\usepackage{amsthm} 
\usepackage{float} 
\usepackage{dcolumn}
\usepackage{titlesec}
\usepackage{enumitem}
\usepackage{hyperref}
\usepackage{tikz}
\usepackage{multicol}
\usetikzlibrary{arrows.meta,positioning,fit,calc}
\tikzset{
  var/.style={rectangle, draw, rounded corners, minimum width=18mm, minimum height=7mm, align=center},
  latent/.style={ellipse, draw, dashed, minimum width=18mm, minimum height=7mm, align=center},
  controlarrow/.style={->, dashed, draw=red!60},
  lightbluearrow/.style={->, dashed, draw=blue!40} % NEW style
}

%------------ Page Style ------------
\rhead{\thepage}
\lhead{ECO1465 Weekly Report 1}

%------------ Section Formatting ------------
\titleformat{\section}{\large\bfseries}{\thesection.}{0.5em}{}
\titleformat{\subsection}{\normalsize\bfseries\itshape}{\thesubsection}{0.5em}{}
\titlespacing*{\subsection}{0pt}{1.5ex plus 1ex minus .2ex}{1em} % spacing before/after subsections


%------------ Begin Document ------------
\begin{document}
% Title Page only
\begin{titlepage}
    \thispagestyle{empty} % no headers/footers
    \centering
    \vspace*{\fill}

    {\Large \textbf{Data Manipulation} \par}

    \vspace{3em}
    {\large Nardeen Abdulkareem \par}
    {\normalsize The University of Toronto \par}
    {\normalsize \today \par}
    {\texttt{nardeen.abdulkareem@mail.utoronto.ca} \par}

    \vspace*{\fill}
\end{titlepage}
\clearpage

\section{Satellite-Based Time Series Diagnostics}

To better understand the behavior of the satellite-derived indicators used in the analysis, we construct monthly time series for several variables derived from the Google Earth Engine extraction. The variables include night-time lights (NTL), night land surface temperature (LST), cloud-free observation coverage, vegetation and built-environment indices (NDVI and NDBI), the Sentinel-2 shortwave infrared band (B11), spectral ratios, and several engineered features intended to capture potential oil and gas activity.

\subsection{Panel Construction}

The raw satellite dataset contains monthly observations for a set of spatial grid cells indexed by $i$ and time periods indexed by $t$. Each observation corresponds to a grid cell of approximately $1 \times 1$ kilometer covering the Kurdish region of Iraq. For each grid cell and month we observe a vector of remotely sensed indicators:

\[
X_{it} = \left( \text{NTL}_{it}, \text{LST}_{it}, \text{NDVI}_{it}, \text{NDBI}_{it}, B11_{it}, \text{CF}_{it}, \ldots \right).
\]

Before constructing the time series, several preprocessing steps are applied:

\begin{enumerate}
    \item Observations coded as $-9999$ are treated as missing values.
    \item Duplicate grid--month observations are collapsed by taking the mean.
    \item Variable-specific transformations are applied where appropriate. In particular,
    \begin{itemize}
        \item $\log(1 + x)$ transformations are applied to night-time lights and cloud-free coverage to reduce skewness.
        \item Log transformations are applied to strictly positive spectral bands such as the SWIR band (B11) and the SWIR/NIR ratio.
    \end{itemize}
\end{enumerate}

The resulting panel therefore contains a cleaned monthly observation for each grid cell and variable.

\subsection{Aggregate Monthly Series}

For each variable $v$, we construct an aggregate regional time series by averaging across all grid cells observed in month $t$. Specifically, the cross-sectional mean is computed as

\[
\bar{X}_{vt} =
\frac{1}{N_t}
\sum_{i=1}^{N_t} X_{vit},
\]

where $N_t$ denotes the number of grid cells with non-missing observations in month $t$.

To provide information about cross-sectional dispersion, we also compute the interquartile range of the distribution across grid cells:

\[
Q_{25,vt}, \qquad Q_{75,vt}.
\]

These statistics are visualized in the plots as a shaded ribbon surrounding the mean time series.

\subsection{Smoothing}

Satellite-derived indicators can exhibit substantial short-run noise due to atmospheric conditions, sensor variation, and compositing artifacts. To reduce high-frequency fluctuations, we apply a centered rolling mean smoother to the monthly series:

\[
\tilde{X}_{vt}
=
\frac{1}{k}
\sum_{j=-k/2}^{k/2} \bar{X}_{v,t+j},
\]

where $k$ denotes the smoothing window (six months in our implementation). The smoothed series is plotted as the main line in each figure, while individual monthly means are displayed as points.

\subsection{Deseasonalization}

Many remote sensing variables exhibit strong seasonal patterns driven by vegetation cycles, surface moisture, and atmospheric conditions. To remove these predictable seasonal components, we compute grid-specific calendar-month averages:

\[
\mu_{vi,m} =
E\left[ X_{vit} \mid \text{month}(t)=m \right].
\]

The deseasonalized series is then defined as

\[
X^{DS}_{vit} =
X_{vit} - \mu_{vi,m}.
\]

This procedure removes systematic month-of-year effects while preserving longer-run deviations that may reflect structural changes in land use or industrial activity.

\subsection{Anomaly Series}

Finally, we construct anomaly series measuring deviations from each grid cell’s typical seasonal level:

\[
A_{vit} =
X_{vit} - \mu_{vi,m}.
\]

These anomaly measures highlight unusual increases or decreases relative to the grid cell’s historical seasonal pattern. Aggregating these anomalies across grid cells provides a regional measure of abnormal satellite signals over time.

\subsection{Interpretation}

The resulting figures therefore present three complementary views of each satellite indicator:

\begin{enumerate}
    \item \textbf{Raw series}: the regional mean level of the variable.
    \item \textbf{Deseasonalized series}: the variable with predictable seasonal effects removed.
    \item \textbf{Anomaly series}: deviations from each grid cell's historical seasonal baseline.
\end{enumerate}

The raw series illustrates broad temporal patterns in the underlying satellite indicators, while the deseasonalized and anomaly series help isolate potential structural changes in activity. This distinction is particularly important for remotely sensed variables such as NDVI or NDBI, which can display strong seasonal fluctuations unrelated to underlying industrial activity.

Together, these diagnostics provide a transparent overview of the behavior of the satellite-derived indicators prior to their use in the supervised learning framework used to identify potential oil and gas extraction sites.










\end{document}